% SOURCE_AUTHORITY: sga5_fr_workpass.tex, lines 580--630 (SHA256 791F4EFFC5E02832D5D77ED03518C8156D6F07E4C8238B03545DB93D883FBB28; LF-normalized unit).
\subsection*{2. Unicidade do complexo dualizante}

Neste número, supomos que $A$ é local, isto é, que $A = \Z/\ell^{\nu}\Z$, com $\ell$ primo.

\subsection*{Teorema 2.1 (cf. [H] V 3.1)}
Seja $X$ um esquema (localmente noetheriano) conexo, e seja $K$ um complexo dualizante sobre $X$ (1.6).

Seja, por outro lado, $K' \in \ob \D_c^{+}(X)$. Para que $K'$ seja dualizante, é necessário e suficiente que existam um $A_X$-módulo inversível $L$ e um inteiro $r$ (que ficam então determinados a menos de isomorfismo único) tais que:
\[
K' \simeq K \mathop{\otimes}\limits^{\mathrm{L}} L[r].
\]

\emph{Demonstração.} Sejam $L$ um $A_X$-módulo inversível e $r$ um inteiro, e ponhamos $L^{\vee} = \underline{\Hom}(L, A_X)$. Existem, para $F \in \ob \D(X)$ e $G \in \ob \D^{+}(X)$, isomorfismos canônicos:
\[
(*) \quad \R\,\underline{\Hom}(F \mathop{\otimes}\limits^{\mathrm{L}} L^{\vee}[-r], G) \simeq \R\,\underline{\Hom}(F, G \mathop{\otimes}\limits^{\mathrm{L}} L[r]) \simeq \R\,\underline{\Hom}(F,G) \mathop{\otimes}\limits^{\mathrm{L}} L[r].
\]
Daí se deduz facilmente que, se $K' \simeq K \mathop{\otimes}\limits^{\mathrm{L}} L[r]$, então $K'$ é dualizante.

Além disso, se colocarmos $\underline{D}_K = \underline{D}$, $\underline{D}_{K'} = \underline{D}'$, temos, para $F \in \ob \D(X)$, $\underline{D}'F \simeq \underline{D} F \mathop{\otimes}\limits^{\mathrm{L}} L[r]$ por $(*)$; portanto, como $\underline{D} K \simeq A_X$, temos
\[
\underline{D}' K \simeq L[r],
\]
O que mostra que $L[r]$, e portanto $L$ e $r$, são determinados de forma essencialmente única por $K'$.

Reciprocamente, suponhamos que $K'$ é dualizante, e ponhamos
\[
P = \underline{D}'K = \underline{D}'\underline{D} A_X = \R\,\underline{\Hom}(K,K').
\]
Vamos demonstrar que se tem
\begin{align*}
\text{a)} \qquad K' &\simeq K \mathop{\otimes}\limits^{\mathrm{L}} P, \\
\text{b)} \qquad P &\simeq L[r], \text{ para um $A_X$-módulo inversível $L$, e assim terminará a demonstração do teorema.}
\end{align*}

Observemos primeiro que, em virtude de 1.9~b)(ii), tem-se, para $F \in \ob \D_c(X)$,
\[
(**) \qquad F \mathop{\otimes}\limits^{\mathrm{L}} P \simeq \underline{D}'\underline{D} F.
\]

Aplicando $(**)$ a $F = K$, obtém-se
\[
K \mathop{\otimes}\limits^{\mathrm{L}} P \simeq \underline{D}'\underline{D} K \simeq \underline{D}' A_X = K', \text{ que demonstra a).}
\]

Aplicando $(**)$ a $F = P' = \underline{D}\underline{D}'A_X$, obtém-se
\[
P' \mathop{\otimes}\limits^{\mathrm{L}} P = \underline{D}'\underline{D}\underline{D}\underline{D}'A_X \simeq \underline{D}'\underline{D}'A_X \simeq A_X.
\]
Portanto, b) será consequência do lema seguinte:
